\section{Phần C --- Bài tập tính toán và phân tích}
\label{sec:phan-c}

Phần~C nối lý thuyết ở Phần~B với mã ở Phần~A: mỗi bài được giải bằng tay hoặc
bằng công thức tường minh, rồi đối chiếu số với cài đặt của nhóm trên CPU
\cite[tr.~19]{de}.

\textbf{Về đầu vào chính thức.} Đề ghi C.1 dùng ``ma trận Q/K/V do giảng viên cung
cấp''. Theo thông tin nhóm nhận được, giảng viên không phát bộ ma trận riêng; C.4
nêu ``cho log-xác suất'' nhưng không kèm số. Nhóm vì vậy dùng \emph{bộ minh họa có
ghi nguồn và tái lập được}, và ghi rõ đây không phải số liệu chính thức. Nếu nhận
được đầu vào riêng, chỉ cần thay số ở ô đầu vào của notebook; toàn bộ phép tính và
phần đối chiếu giữ nguyên.

%--------------------------------------------------------------------
\subsection{C.1 --- Tính tay một đầu attention có mặt nạ nhân quả}
\label{sec:c1}

\textbf{Đề bài.} Ba token, một đầu chú ý, $d_k=2$, có mặt nạ nhân quả.

\textbf{Đầu vào minh họa.} Sinh bằng \texttt{random.Random(577893521)}, lấy số
nguyên trong $[-2,2]$ theo hàng, lần lượt $Q$, $K$, $V$; không lấy mẫu lại. Mỗi
hàng là một token; chọn $d_v=2$.
\[
  Q=\begin{bmatrix}-2&0\\-1&-2\\1&-2\end{bmatrix},\quad
  K=\begin{bmatrix}-1&-1\\2&1\\0&2\end{bmatrix},\quad
  V=\begin{bmatrix}-1&0\\-1&1\\0&-1\end{bmatrix}.
\]

\subsubsection{Bước 1: điểm số, chia tỷ lệ, che tương lai}

Phần tử $(i,j)$ của $QK^\top$ là tích vô hướng của query token $i$ với key token
$j$. Ví dụ hàng 2: $(-1)(-1)+(-2)(-1)=3$; $(-1)2+(-2)1=-4$; $(-1)0+(-2)2=-4$.
Đặt $c=1/\sqrt{2}$:
\[
  QK^\top=\begin{bmatrix}2&-4&0\\3&-4&-4\\1&0&-4\end{bmatrix},
  \qquad
  S=c\,QK^\top+M=\begin{bmatrix}
      2c&-\infty&-\infty\\
      3c&-4c&-\infty\\
      c&0&-4c\end{bmatrix}.
\]
Mặt nạ cộng $M_{ij}=0$ nếu $j\le i$ và $-\infty$ nếu $j>i$. Việc che xảy ra
\emph{trước} khi chuẩn hóa từng hàng: vì $e^{-\infty}=0$, token tương lai không góp
vào cả tử \emph{và} mẫu. Đường chéo được giữ, nên mỗi hàng luôn còn ít nhất một
khóa hợp lệ.

\subsubsection{Bước 2: softmax từng hàng và tổng có trọng số}

Gọi $A=\softmax_{\mathrm{hàng}}(S)$ và $o_i=\sum_{j\le i}A_{ij}v_j$. Trừ điểm lớn
nhất của hàng trước khi lấy mũ không làm đổi softmax (tính bất biến dịch), và
tránh tràn số.

\begin{description}[leftmargin=2em,style=nextline]
\item[Token 1] Chỉ $v_1$ hợp lệ, nên $A_{1,:}=(1,0,0)$ và $o_1=v_1=(-1,0)$.
\item[Token 2] Trừ $3c$ khỏi hàng 2. Đặt $t=e^{-7c}$, khi đó
$A_{2,:}=(1,t,0)/(1+t)$ và
$o_2=\bigl[(-1,0)+t(-1,1)\bigr]/(1+t)=\bigl(-1,\ t/(1+t)\bigr)$.
Tọa độ đầu vẫn đúng bằng $-1$ vì cả hai value hợp lệ đều có tọa độ đầu bằng $-1$.
\item[Token 3] Trừ $c$ khỏi hàng 3. Đặt $u=e^{-c}$, $v=e^{-5c}$, $D=1+u+v$. Khi đó
$A_{3,:}=(1,u,v)/D$ và
$o_3=\bigl[(-1,0)+u(-1,1)+v(0,-1)\bigr]/D=\bigl(-(1+u)/D,\ (u-v)/D\bigr)$.
\end{description}

\begin{center}
\begin{tabular}{crr}
\toprule
\textbf{Token} & \textbf{Trọng số trên key 1, 2, 3} & \textbf{Đầu ra hai chiều}\\
\midrule
1 & $(1,\ 0,\ 0)$ & $(-1,\ 0)$\\
2 & $(0{,}992965,\ 0{,}007035,\ 0)$ & $(-1,\ 0{,}007035)$\\
3 & $(0{,}656939,\ 0{,}323916,\ 0{,}019145)$ & $(-0{,}980855,\ 0{,}304771)$\\
\bottomrule
\end{tabular}
\end{center}

Bảng được làm tròn \emph{sau} khi tính $O=AV$; không dùng trọng số đã làm tròn để
tính tiếp.

\subsubsection{Đối chiếu \dalam}

Notebook C.1 chạy \texttt{scaled\_dot\_product\_attention} và \texttt{causal\_mask}
của A.1 ở \texttt{float64} trên CPU, so với biểu thức đóng ở trên bằng
\texttt{torch.testing.assert\_close} với \texttt{atol=1e-12}, và kiểm thêm
\texttt{weights.triu(1).count\_nonzero() == 0} --- tức không có trọng số dương nào
ở phần tam giác trên, xác nhận không rò rỉ tương lai.

\textbf{Nhận xét đáng nêu ở vấn đáp.} Trọng số của token 2 đặt $99{,}3\%$ khối
lượng vào key 1 vì chênh lệch điểm là $7c\approx4{,}95$ --- đúng hiện tượng bão hòa
mô tả ở \eqref{eq:b1-sat} với $m=2$, $g=7c$. Đây là ví dụ cụ thể cho thấy vì sao
hệ số $1/\sqrt{d_k}$ có ý nghĩa.

%--------------------------------------------------------------------
\subsection{C.2 --- Đếm tham số và phép tính của một khối Transformer}
\label{sec:c2}

\textbf{Phạm vi.} Một lần forward của \emph{một khối} trên một chuỗi
$X\in\R^{n\times d}$, với $d=\dmodel$, $f=d_{ff}$, $h$ đầu, $h\mid d$,
$d_h=d/h$. Không tính embedding, PE, LayerNorm cuối mô hình, đầu đọc từ vựng,
backward hay optimizer. Mặt nạ nhân quả là buffer cố định, không phải tham số.

\subsubsection{Số tham số}

\texttt{Linear(a,b)} học $ab$ trọng số và $b$ bias. LayerNorm trên $d$ chiều học cả
$\gamma$ và $\beta$, mỗi vector dài $d$. Phép reshape chia đầu chỉ tổ chức lại cùng
$d$ chiều, \emph{không} nhân số tham số lên $h$ lần.

\begin{center}
\begin{tabular}{lrrr}
\toprule
\textbf{Thành phần} & \textbf{Trọng số} & \textbf{Bias} & \textbf{Tổng}\\
\midrule
Chiếu Q,K,V gộp ($d\to3d$) & $3d^2$ & $3d$ & $3d^2+3d$\\
Chiếu đầu ra ($d\to d$) & $d^2$ & $d$ & $d^2+d$\\
FFN lớp 1 ($d\to f$) & $df$ & $f$ & $df+f$\\
FFN lớp 2 ($f\to d$) & $df$ & $d$ & $df+d$\\
Hai LayerNorm & $2d$ & $2d$ & $4d$\\
\midrule
\multicolumn{3}{r}{\textbf{Tổng}} & $\mathbf{4d^2+2df+9d+f}$\\
\bottomrule
\end{tabular}
\end{center}

\begin{equation}
  N_{\mathrm{params}}=4d^2+2df+9d+f.
  \label{eq:c2-params}
\end{equation}

Với $d,f$ cố định, số tham số \emph{không} phụ thuộc $n$ hay $h$. ReLU, softmax,
phép cộng residual và mặt nạ không có tham số học.

\subsubsection{Số phép tính forward}

Quy ước MAC: tích $(a\times b)(b\times c)$ dùng $abc$ lượt nhân--tích lũy; khi đối
chiếu profiler dùng 1 MAC $=$ 2 FLOP theo quy ước GEMM.

\begin{center}
\begin{tabular}{llr}
\toprule
\textbf{Thành phần} & \textbf{Kích thước tích} & \textbf{MAC cả khối}\\
\midrule
Chiếu Q,K,V & $(n\times d)(d\times 3d)$ & $3nd^2$\\
Điểm $QK^\top$ & mỗi đầu $(n\times d_h)(d_h\times n)$ & $hn^2d_h=n^2d$\\
Tổng có trọng số $PV$ & mỗi đầu $(n\times n)(n\times d_h)$ & $hn^2d_h=n^2d$\\
Chiếu đầu ra & $(n\times d)(d\times d)$ & $nd^2$\\
FFN hai lớp & $(n\times d)(d\times f)$, $(n\times f)(f\times d)$ & $2ndf$\\
\midrule
\multicolumn{2}{r}{\textbf{Tổng}} & $\mathbf{4nd^2+2ndf+2n^2d}$\\
\bottomrule
\end{tabular}
\end{center}

\begin{equation}
  \mathrm{MAC}_{\mathrm{matmul}}=4nd^2+2ndf+2n^2d.
  \label{eq:c2-mac}
\end{equation}

Lưu ý phép triệt tiêu $h\,d_h=d$: chi phí attention \emph{không} phụ thuộc số đầu
khi $d$ giữ nguyên. Đây là điểm bị hỏi thường xuyên.

\textbf{Ngoài các tích ma trận:} cộng bias $n(5d+f)$; hai residual $2nd$; ReLU
$\Theta(nf)$; hai LayerNorm $\Theta(nd)$; chia tỷ lệ điểm $hn^2$ phép nhân và
softmax xử lý $hn^2$ phần tử, chi phí $\Theta(hn^2)$.

\subsubsection{Hạng tử bậc hai theo $n$}

Trong \eqref{eq:c2-mac}, hai hạng $n^2d$ đến từ $QK^\top$ và $PV$; phần chiếu và
FFN tăng \emph{tuyến tính} theo $n$. Vì $h\le d$, độ phức tạp toàn khối là
$\Theta(nd^2+ndf+n^2d)$, và score cùng xác suất attention có kích thước
$h\times n\times n$ --- tức bộ nhớ cũng bậc hai. Với batch $B$, số phép tính nhân
$B$ còn số tham số giữ nguyên.

\textbf{Vì sao mặt nạ nhân quả không chia đôi chi phí?} Cài đặt A.1 tính đủ ma trận
$n\times n$ rồi mới điền $-\infty$ phía trên đường chéo, và $PV$ vẫn là phép nhân
dense đủ kích thước dù các trọng số bị che bằng $0$. Chỉ thuật toán \emph{thực sự
bỏ qua} ô bị che mới khai thác được $n(n+1)/2$ cặp hợp lệ; mặt nạ trong A.1 không
làm việc đó. Đây là khác biệt giữa đúng về ngữ nghĩa và tối ưu về chi phí.

\subsubsection{Thay số theo cấu hình A.2 \dalam}

Với $d=64$, $f=128$, $h=4$, $n=128$:
\begin{align*}
  N_{\mathrm{params}}&=4(64)^2+2(64)(128)+9(64)+128\\
   &=16{.}384+16{.}384+576+128=\mathbf{33{.}472},\\
  \mathrm{MAC}&=4(128)(64)^2+2(128)(64)(128)+2(128)^2(64)\\
   &=2{.}097{.}152+2{.}097{.}152+2{.}097{.}152=\mathbf{6{.}291{.}456}.
\end{align*}
Tương đương $12{.}582{.}912$ FLOP theo quy ước trên. Ở cấu hình này hai tích
attention chiếm \emph{đúng} $1/3$ tổng MAC; tỷ lệ đó tăng khi chuỗi dài hơn.

Notebook C.2 đối chiếu hai chiều: \texttt{numel()} của từng submodule khớp
\eqref{eq:c2-params} cho cả $h=4$ và $h=8$ (xác nhận tính độc lập với $h$), và FLOP
mà PyTorch Profiler ước lượng cho \texttt{aten::addmm} (chiếu, FFN) cùng
\texttt{aten::bmm} (hai tích attention) khớp đúng $2\times$ \eqref{eq:c2-mac} với
$n\in\{64,128,256\}$. Khi $n$ gấp đôi, chiếu và FFN gấp đôi còn hai tích attention
gấp bốn --- xác nhận hạng tử $O(n^2)$ bằng số.

Không dùng số FLOP của một khối forward làm FLOP huấn luyện toàn mô hình.

%--------------------------------------------------------------------
\subsection{C.3 --- Tham số LoRA ở Q/V và bộ nhớ QLoRA 4 bit}
\label{sec:c3}

\textbf{Mô hình nền dùng để tính.} Qwen2.5-1.5B-Instruct, khóa theo revision
\texttt{989aa7980e4cf806f80c7fef2b1adb7bc71aa306} \cite{qwen}. Đây là mô hình dùng
cho \emph{phép tính} C.3; nó không đồng nghĩa với quyết định mô hình của D.3.

Kiến trúc: $L=28$ khối, chiều ẩn $d=1536$, 12 đầu Q nhưng chỉ 2 đầu K/V
(grouped-query attention), mỗi đầu 128 chiều. Do đó trọng số Q có dạng
$1536\times1536$ còn V là $256\times1536$ (đầu ra $\times$ đầu vào). Việc lặp K/V
khi tính attention \emph{không} tạo thêm trọng số.

\subsubsection{Đếm tham số học}

Theo \cite[§4.1]{lora}, với $W$ kích thước $o\times i$, hai ma trận học có kích
thước $r\times i$ và $o\times r$, tức $r(i+o)$ tham số --- đúng
\eqref{eq:b4-count}. Đóng băng toàn bộ mô hình nền kể cả bias; chỉ gắn adapter vào
\texttt{q\_proj} và \texttt{v\_proj} ở mỗi khối, không học K, O, embedding hay đầu
đọc. Hệ số $\alpha/r$ không làm đổi số tham số.
\begin{align*}
  P_Q&=28(1536+1536)r=86{.}016\,r,\\
  P_V&=28(1536+256)r=50{.}176\,r,\\
  P&=136{.}192\,r.
\end{align*}

\begin{center}
\begin{tabular}{rrrr}
\toprule
$r$ & \textbf{Q} & \textbf{V} & \textbf{Tổng}\\
\midrule
4 & 344.064 & 200.704 & \textbf{544.768}\\
8 & 688.128 & 401.408 & \textbf{1.089.536}\\
16 & 1.376.256 & 802.816 & \textbf{2.179.072}\\
\bottomrule
\end{tabular}
\end{center}

Với $r=8$, tỷ lệ tham số học so với $1{,}54$ tỷ tham số của mô hình nền là khoảng
$0{,}071\%$.

Notebook C.3 kiểm số này bằng chính lớp \texttt{LoRALinear} của A.3 trên
\emph{meta tensor} --- tức kiểm kích thước mà không cấp vùng nhớ, không lượng tử
hóa và không tải Qwen.

\subsubsection{Ước lượng bộ nhớ khi kết hợp QLoRA 4 bit}

\textbf{Kịch bản:} SFT, không offload, không có reference policy của DPO. Phân biệt
rõ \emph{trọng số lưu} NF4 với \emph{kiểu tính toán} 16 bit.

Bảy ma trận lớn mỗi khối (Q/K/V/O và gate/up/down) lưu NF4 với double
quantization. Gọi $k=256$ là chiều ra K/V và $f=8960$ là chiều trung gian. Số phần
tử lượng tử hóa là
\[
  N_4=L\,(2d^2+2dk+3df).
\]
Theo \cite[§3, tr.~5]{qlora}, mỗi phần tử dùng trung bình
$4+\tfrac{8}{64}+\tfrac{32}{64\cdot256}$ bit, kể cả scale hai cấp. Bỏ qua codebook,
offset và padding theo tensor.

Giả sử embedding và đầu đọc dùng chung trọng số
(\texttt{tie\_word\_embeddings=true}), RMSNorm và bias lưu 16 bit:
\[
  N_{16}=V_{\mathrm{ocab}}\,d+(2L+1)d+L(d+2k),
\]
gồm một embedding, hai norm mỗi khối, một norm cuối, và bias Q/K/V. Không đếm
\texttt{lm\_head} lần thứ hai. Đây là \emph{giả định lưu trữ}, không phải mặc định
của mọi thư viện: nếu chuyển phần này sang FP32 thì cộng thêm $2N_{16}$ byte.

Adapter lưu FP32, gradient FP32, hai moment AdamW FP32 thường trú, tức
\textbf{16 byte cho mỗi tham số học}; không tính gradient hay optimizer cho
backbone đóng băng. Sau khi trạng thái optimizer đã cấp phát và gradient còn sống:
\begin{equation}
  M_{\mathrm{tĩnh}}=\frac{N_4\,b_4}{8}+2N_{16}+16P \quad\text{byte}.
  \label{eq:c3-static}
\end{equation}

Notebook C.3 có một đối chứng độc lập đáng chú ý: $N_4+N_{16}=1{.}543{.}714{.}304$,
khớp đúng tổng số phần tử BF16 của checkpoint trên Hub trước khi lượng tử hóa. Đây
là phép kiểm bắt được lỗi đếm thiếu hoặc đếm trùng một họ ma trận.

\subsubsection{Từ bộ nhớ tĩnh tới kịch bản có activation}

Chọn microbatch $B=1$, chiều dài $T=512$, checkpoint từng khối, không KV cache.
Ước lượng các tensor chính ở pha tính loss: một đầu vào FP16 cho mỗi khối
($2LBTd$ byte), cùng logits FP32 và một tensor FP32 cùng cỡ cho loss không fused
($2\times4BTV_{\mathrm{ocab}}$ byte).

\textbf{Kết luận định tính.} Tăng rank từ 4 lên 16 làm số tham số học tăng 4 lần
nhưng chỉ thêm khoảng 25 MiB trạng thái adapter theo \eqref{eq:c3-static}, không
nhân bốn bộ nhớ cả mô hình. Phần chi phối ngân sách là backbone lượng tử hóa và các
tensor từ vựng, không phải adapter.

\textbf{Giới hạn \chualam.} Đây là \emph{ước lượng theo các tensor được nêu}, không
phải phép đo. Chưa cộng: tái tính activation và backward của block, buffer giải
lượng tử, workspace GEMM/attention, CUDA context, tensor tạm của optimizer, và phân
mảnh allocator. Khi $T$ tăng, hai khoản activation tăng tuyến tính; nếu attention
lưu điểm dense thì riêng một ma trận điểm FP16 cỡ $B\times h\times T\times T$ tăng
\emph{bậc hai}. Vì vậy \textbf{không} được suy từ bảng rằng một GPU cụ thể chạy vừa
SFT hay DPO. Muốn xác nhận VRAM cho D.3 phải đo cấu hình thật qua ít nhất một bước
forward, backward và optimizer sau khi cấp phát trạng thái, ghi cả peak allocated
và peak reserved.

%--------------------------------------------------------------------
\subsection{C.4 --- Mất mát DPO, gradient và ảnh hưởng của $\beta$}
\label{sec:c4}

\subsubsection{Lời giải cho một cặp bất kỳ}

Cùng một prompt $x$, gọi $w$ là đáp án được ưa thích, $l$ là đáp án bị từ chối.
Đặt $p_i=\log\pi_\theta(y_i\mid x)$ và $r_i=\log\pref(y_i\mid x)$ với
$i\in\{w,l\}$. Đây là \emph{tổng} log-xác suất tự nhiên của token đáp án, không gồm
prompt hay padding, và không phải trung bình theo độ dài như trong perplexity
(Mục~\ref{sec:b6-agg}). Các log-xác suất hữu hạn; $\beta>0$; reference đóng băng.

Theo \eqref{eq:b5-logit} và \cite[§4, Eq. (7)]{dpo}:
\begin{equation}
  m=(p_w-p_l)-(r_w-r_l),\qquad z=\beta m,\qquad
  L=-\log\sigma(z)=\log(1+e^{-z}).
  \label{eq:c4-loss}
\end{equation}
Điểm cần nhấn: $m$ so \emph{tỷ số} xác suất của hai đáp án giữa policy và
reference, không chỉ so $p_w$ với $p_l$. Và $\sigma(z)$ là xác suất \emph{ưu tiên},
khác xác suất policy sinh ra riêng đáp án $w$.

Vì $\mathrm dL/\mathrm dz=-\sigma(-z)$ và
$\partial z/\partial p_w=\beta=-\partial z/\partial p_l$, đặt $g=\beta\sigma(-z)$:
\begin{equation}
  \frac{\partial L}{\partial p_w}=-g,\qquad
  \frac{\partial L}{\partial p_l}=+g,\qquad
  \frac{\partial L}{\partial r_w}=\frac{\partial L}{\partial r_l}=0 .
  \label{eq:c4-grad}
\end{equation}
Hai gradient đối nhau: mất mát đẩy log-xác suất của $w$ lên và của $l$ xuống với
cùng độ lớn. Gradient theo reference bằng $0$ vì reference bị
\texttt{detach} --- phép kiểm trong notebook xác nhận điều này bằng
\texttt{assert grw is None and grl is None}.

\subsubsection{Thế số --- ví dụ minh họa}

Chọn $(p_w,p_l,r_w,r_l)=(\log 0{,}2,\ \log 0{,}1,\ \log 0{,}4,\ \log 0{,}1)$ và
$\beta=1$. Hai đáp án không phải toàn bộ không gian đầu ra, nên hai xác suất của
mỗi mô hình không cần cộng thành $1$.

Policy ưu tiên $w$ so với $l$ theo tỷ số $2:1$, nhưng reference ưu tiên mạnh hơn,
$4:1$. Vậy
\[
  m=\log 2-\log 4=-\log 2,\quad z=-\log 2,\quad \sigma(z)=\tfrac13,
\]
\[
  \boxed{L=\log 3\approx 1{,}098612},\qquad
  g=\tfrac23,\qquad
  \left(\frac{\partial L}{\partial p_w},\frac{\partial L}{\partial p_l}\right)
   =\left(-\tfrac23,\ +\tfrac23\right).
\]

\textbf{Quan sát quan trọng:} mất mát $\log 3$ \emph{lớn hơn} $\log 2$ dù
$p_w>p_l$. Lý do: policy đã ưu tiên $w$, nhưng chưa ưu tiên mạnh \emph{bằng}
reference, nên $m<0$. Đây là điểm dễ bị hỏi vặn nhất ở C.4, và nó chính là hệ quả
của việc \eqref{eq:b5-logit} so hai log-tỷ số chứ không so hai log-xác suất.

\subsubsection{Ảnh hưởng của $\beta$ tại một cặp cố định}

Ở một cặp cố định, $\partial L/\partial\beta=-m\,\sigma(-\beta m)$. Khi
$\beta\to0^+$, mất mát tiến tới $\log 2$ và $g\sim\beta/2$: tín hiệu gradient yếu
cho \emph{mọi} cặp.

\begin{center}
\begin{tabular}{p{0.32\textwidth}p{0.28\textwidth}p{0.26\textwidth}}
\toprule
\textbf{Policy so với reference} & \textbf{Khi tăng $\beta$} & \textbf{Khi $\beta$ rất lớn}\\
\midrule
$m<0$: chênh lệch chưa đủ & Mất mát tăng; $|g|$ tăng & $L\sim\beta|m|$, $g\sim\beta$\\
$m=0$: hai tỷ số bằng nhau & $L=\log2$ không đổi, nhưng gradient \emph{không} bằng $0$ & $g=\beta/2$\\
$m>0$: đã vượt reference & Mất mát giảm; $g$ không tăng đơn điệu & $L\to0$, $g\to0$ do bão hòa\\
\bottomrule
\end{tabular}
\end{center}

Hàng giữa đáng chú ý: khi hai tỷ số bằng nhau, mất mát bằng $\log 2$ \emph{bất kể}
$\beta$, nhưng gradient bằng $\beta/2\ne0$ --- mô hình vẫn bị đẩy đi. Đây là khác
biệt giữa ``mất mát không đổi'' và ``không học gì''.

\textbf{Phân biệt hai cách đọc $\beta$.} Bảng trên chỉ đổi $\beta$ tại policy cố
định. Trong mục tiêu gốc $\E[R]-\beta\KL(\pi\|\pref)$ ở \eqref{eq:b5-obj}, giữ cố
định reward và reference thì $\beta$ lớn đặt trọng số cao hơn lên phạt KL. Với cùng
prompt và hai đáp án có xác suất reference dương, \eqref{eq:b5-gibbs} cho tỷ số tối
ưu
\[
  \frac{\pi^*(y_w\mid x)}{\pi^*(y_l\mid x)}
   =\frac{\pref(y_w\mid x)}{\pref(y_l\mid x)}
    \exp\!\left(\frac{R_w-R_l}{\beta}\right).
\]
Nghiệm thay đổi theo $\beta$ \emph{nếu} reward không hằng trên hỗ trợ của reference.
Nếu $R(x,y)=c(x)$ trên hỗ trợ đó thì hệ số $e^{c(x)/\beta}$ là chung cho mọi đáp án,
bị triệt tiêu khi chuẩn hóa, và $\pi^*=\pref$ với mọi $\beta>0$. Hai cách đọc này
không được lẫn.

\subsubsection{Đối chiếu \dalam}

Notebook C.4 kiểm trên CPU ở \texttt{float64}, dùng hai dạng ổn định số cho
$\mathrm{softplus}(-z)$ và $\sigma(-z)$, so với \texttt{src/dpo.py} và với
\texttt{torch.autograd.grad}: mất mát và cả hai gradient policy khớp ở
\texttt{rel\_tol=1e-12}; gradient theo reference xác nhận là \texttt{None}. Phép
kiểm chạy qua \textbf{ba dấu của margin} ($m=-\log2,\ 0,\ +\log2$, bằng cách đổi
xác suất policy cho $w$ thành $0{,}2$, $0{,}4$, $0{,}8$) và \textbf{bốn giá trị}
$\beta\in\{0{,}1;\ 1;\ 2;\ 8\}$. Mỗi hàng là một cặp riêng; không lấy trung bình
giữa các ví dụ.

\textbf{Giới hạn.} Bộ số là minh họa, không phải số giảng viên cấp, không phải đầu
ra của một checkpoint, và không phải nhãn người chấm ở D.4. Không dùng bảng này để
chọn $\beta$ cho D.4 hay kết luận DPO cải thiện chất lượng.

%--------------------------------------------------------------------
\subsection{Tổng hợp trạng thái Phần C}

\begin{center}
\begin{tabular}{llp{0.42\textwidth}}
\toprule
\textbf{Mã} & \textbf{Trạng thái} & \textbf{Kết quả chính}\\
\midrule
C.1 & \dalam & Ba token giải từng hàng, khớp A.1 ở \texttt{atol=1e-12}; đầu vào minh họa seed \texttt{577893521}\\
C.2 & \dalam & 33.472 tham số, 6.291.456 MAC tại $n=128$; hạng tử $O(n^2)$ xác nhận bằng profiler\\
C.3 & \dalam & 544.768 / 1.089.536 / 2.179.072 tham số học; ước lượng QLoRA, chưa đo VRAM\\
C.4 & \dalam & $L=\log3$, gradient $(-2/3,+2/3)$; khớp A.4 và autograd trên 3 dấu margin $\times$ 4 giá trị $\beta$\\
\bottomrule
\end{tabular}
\end{center}
